\chapter{Rendszerszintű értékelés és megvitatás}

\section{A legfontosabb rendszertervezési döntések}

Az egész projektet visszanézve a legerősebb döntésnek az tekinthető, hogy a különböző feladatok nem kerültek egyetlen univerzális technológiába kényszerítve. A rendszer nem eszközpreferencia, hanem minőségi attribútumok mentén fejlődött: reprodukálhatóság, magyarázhatóság, perzisztencia, futási teljesítmény és demo-koherencia. Ezeket ritkán lehet mindent ugyanazzal az eszközzel a legjobban kiszolgálni.

A Python gyors iterációt és kényelmes adatfeldolgozást ad, a Rust teljesítményt és pontos gráfkontrollt, a Node/Express stabil integrációs felületet, a React pedig erős bemutathatóságot. A többnyelvűség tehát nem öncélú komplexitás, hanem funkcionális specializáció.

\section{A rendszer erősségei}

A projekt legfontosabb erőssége, hogy a teljes adatút végigkövethető a Riot API-tól a frontendig: nincsenek kézzel összeollózott közbenső lépések. A játékosprofilok explicit, ellenőrizhető képleteken és mezőkön alapulnak; a gráfépítés és klaszterezés nem csak vizuális export, hanem perzisztens objektum, amelynek bejárása és újraelőállítása determinisztikus. A Rust réteg valódi analitikai és futtatási értéket ad, nem csupán egy korábbi Python-szkript gyorsabb átírása. Az asszortativitás, a core-periphery gráfolvasat és a Brandes-centralitás valódi kutatási kérdésekké emelik a rendszert, amelyeket az implementáció közvetlenül tud megválaszolni.

\section{Érvényességi fenyegetések}

A projekt legfontosabb validitási kockázata, hogy a dataset nem reprezentálja a teljes játékospopulációt: a gyűjtés seed- és queue-függő, ezért a mintavétel nem véletlen. Az \textit{opscore} és \textit{feedscore} értelmezhető, de nem objektív igazságmérő: különösen az \textit{opscore} tartalmaz kézi szerepkör-súlyozást. A retired signed-balance ág nem kap fő empirikus szerepet, mert az enemy-él negatív társas kötésként való értelmezése nem elég erős. Az ország-elemzési ág sokkal bizonytalanabb, mint a közvetlen gráf- és meccsanalitika, és ezért nem képezi a főeredmények részét.

Rétegenként szétbontva: a \textbf{mintavételi validitás} szempontjából a seedválasztás, a premade-probe és a queue-orientáció együtt nem véletlen mintát adnak a teljes játékospopulációból. A \textbf{mérési validitás} szempontjából az \textit{opscore} és \textit{feedscore} értelmezhetők, de konstrukciós döntésekre épülnek. A \textbf{projekciós validitás} szempontjából a gráfmód, az élküszöb és az ally/enemy kapcsolatértelmezés érdemben alakítja a végső analitikai mintát: a signed-balance emiatt csak módszertani határesetként kezelhető. A \textbf{klasztervaliditás} szempontjából a közösségdetektálás eredménye függ a választott gráftól, a küszöböléstől és a modularitás felbontási korlátjától \cite{fortunato_barthelemy2007}. A \textbf{temporális validitás} szempontjából a korpusz fokozatosan épülő archívum, ezért a teljes háló nem egyetlen szinkron időpillanat felvétele.

Ezek a fenyegetések nem rombolják le a dolgozatot. Inkább kijelölik azt a határt, amelyen belül az állítások védhetők, és a bírálat szempontjából sokkal erősebb pozíciót ad egy rétegzetten dokumentált korlát, mint egy rejtett módszertani gyengeség.

\section{A megvalósult és a jövőbeli elemek szétválasztása}

A dolgozat egyik tudatos módszertani döntése az volt, hogy a jelenleg implementált és a még csak tervezett elemeket nem keveri össze. A ténylegesen megvalósult mag a Riot API-ra épülő adatgyűjtésből, \textit{puuid}-alapú normalizálásból, \textit{feedscore}/\textit{opscore} számításból, legacy Python exportból és populációs közösségdetektálásból, SQLite klaszterperzisztenciából, Rust Graph V2 gráfépítésből, útvonalkeresésből, asszortativitásból, Brandes-centralitásból és 3D globális exportból áll. Ehhez adódik a databaseproject és C\# böngésző alprojekt, amely formálisan nem a főrendszer része, de valódi, futtatható eszközként él a repositoryban.

A temporal consistency és community cohesion vs. performance vizsgálatok eltávolítva a jelenlegi scope-ból. A Brandes-centralitás ezzel szemben már implementált modul, amelynek további feladata nem az alapalgoritmus megírása, hanem a benchmarkolás és a Graph V2 artifact-rétegbe illesztés.

Ez a szétválasztás nem gyengíti, hanem erősíti a dolgozatot. A bírálat szempontjából sokkal védhetőbb egy olyan szöveg, amely pontosan kijelöli a rendszer jelenlegi határát, mint egy olyan, amely minden jó ötletet kész eredménynek próbál feltüntetni.
